% Options for packages loaded elsewhere
\PassOptionsToPackage{unicode}{hyperref}
\PassOptionsToPackage{hyphens}{url}
\documentclass[smallextended]{svjour3}
\usepackage{xcolor}
\usepackage{amsmath,amssymb}
\usepackage{iftex}
\ifPDFTeX
  \usepackage[T1]{fontenc}
  \usepackage[utf8]{inputenc}
  \usepackage{textcomp} % provide euro and other symbols
\else % if luatex or xetex
  \usepackage{unicode-math} % this also loads fontspec
  \defaultfontfeatures{Scale=MatchLowercase}
  \defaultfontfeatures[\rmfamily]{Ligatures=TeX,Scale=1}
\fi
\usepackage{lmodern}
\ifPDFTeX\else
  % xetex/luatex font selection
\fi
% Use upquote if available, for straight quotes in verbatim environments
\IfFileExists{upquote.sty}{\usepackage{upquote}}{}
\IfFileExists{microtype.sty}{% use microtype if available
  \usepackage[]{microtype}
  \UseMicrotypeSet[protrusion]{basicmath} % disable protrusion for tt fonts
}{}
\usepackage{longtable,booktabs,array}
\newcounter{none} % for unnumbered tables
\usepackage{calc} % for calculating minipage widths
% Correct order of tables after \paragraph or \subparagraph
\usepackage{etoolbox}
\makeatletter
\patchcmd\longtable{\par}{\if@noskipsec\mbox{}\fi\par}{}{}
\makeatother
% Allow footnotes in longtable head/foot
\IfFileExists{footnotehyper.sty}{\usepackage{footnotehyper}}{\usepackage{footnote}}
\makesavenoteenv{longtable}
\setlength{\emergencystretch}{3em} % prevent overfull lines
\providecommand{\tightlist}{%
  \setlength{\itemsep}{0pt}\setlength{\parskip}{0pt}}
\usepackage{bookmark}
\IfFileExists{xurl.sty}{\usepackage{xurl}}{} % add URL line breaks if available
\urlstyle{same}
\hypersetup{
  pdftitle={Cycle pairs around Erdős 585: certified extraction{,} structural reductions and recolouring barriers},
  pdfauthor={Robert Huynh; Alejandro Zarzuelo Urdiales},
  hidelinks,
  pdfcreator={LaTeX via pandoc}}

\title{Cycle pairs around Erdős 585: certified extraction, structural reductions and recolouring barriers}
\titlerunning{Cycle pairs around Erdős 585}
\author{Robert Huynh \and Alejandro Zarzuelo Urdiales}
\authorrunning{R. Huynh and A. Zarzuelo Urdiales}
\institute{Robert Huynh (corresponding author) \at \href{https://www.hbs.edu/}{Harvard Business School} \and Alejandro Zarzuelo Urdiales \at \href{https://rsihouse.ai/openmath}{Open Math}}
\date{} % Journal received/accepted dates are supplied by editors.

\begin{document}
\maketitle

\begin{abstract}
We study two edge-disjoint cycles having the same vertex set in finite
simple graphs. For bipartite graphs of maximum degree six, we prove
extraction of a nonempty four-regular subgraph at the threshold
\(3n-5\). Repeated neighbourhoods give a cycle-pair criterion in the
dense range, together with an ordinary six-regular extension allowing
two singleton neighbourhood classes. We give explicit pair-free
five-regular graphs on 18 and 32 vertices and a bipartite example on 104
vertices, and describe finite extremal profiles and incidence-based
lower constructions. A gadget-substitution theorem preserves
pair-freeness at degrees at most seven and yields an equivalence with
prescribed-vertex deletion. Four colours force a pair in prime-field
Haar graphs, with Hamilton decomposition of every exact-four-colour
component; a specified cyclic composite family is also treated. A
canonical parabola colouring defeats every one-round component-union
recolouring detector, while two rounds expose a pair on 54 vertices. We
state and certify quantitative sampling and regular-factor lemmas used
in a separate Hamiltonicity programme. Human-readable constructions and
proof arguments accompany pinned Lean sources and finite certificates.
The result catalogue distinguishes proved statements, conditional
consequences and source-reported computations; no general asymptotic
improvement for the parent extremal problem is claimed.
\keywords{edge-disjoint cycles \and regular subgraphs \and graph substitution \and Haar graphs \and recolouring barriers}
\subclass{Primary 05C38 \and Secondary 05C35, 05C25, 05C70}
\end{abstract}

\begin{center}\small\href{https://github.com/alejandrozu/openmath-2026-judging/tree/70c5ec3107658182f22c207d4c50f53f9cb3e602/robert-review-2026-10-07}{Code and formal proofs}\end{center}

\section{The problem and the
contributions}\label{the-problem-and-the-contributions}

Throughout, a \emph{pair} in a simple graph \(G\) means two genuine
connected cycles \(C_1,C_2\) with \(V(C_1)=V(C_2)\) and
\(E(C_1)\cap E(C_2)=\varnothing\). A graph is \emph{pair-free} if it
contains no pair. Write \(f(n)\) for the maximum number of edges in a
pair-free simple graph on \(n\) vertices. The problem is to determine
the growth of \(f(n)\).

The established general bounds are
\(\Omega(n\log\log n)\leq f(n)\leq n(\log n)^{O(1)}\). The lower-bound
construction of Pyber, Rödl and Szemerédi contains no four-regular
subgraph; the polylogarithmic upper bound is due to Chakraborti, Janzer,
Methuku and Montgomery \cite{ref3,ref2}. Our results concern finite cases,
restricted graph classes, intermediate extraction and the limitations of
a particular proof method. We claim no improvement in that general
asymptotic order.

A pair has a four-regular union on its common support. The converse
fails: a four-regular graph can decompose into two disconnected
two-factors rather than two connected cycles. This distinction separates
the quartic-extraction theorem from the cycle-pair forcing results
throughout the paper.

The eight families are mathematical groupings, not eight independent
solutions of the parent problem. We consolidate weaker variants,
repeated examples, support lemmas and immediate consequences within
their family. Historical novelty is stated relative to the primary
sources actually compared, rather than inferred from formal verification
or from the number of declarations.

Table~\ref{tab:family-map} summarizes the eight result families and their evidence boundaries.

{
\begin{longtable}[]{@{}
  >{\raggedright\arraybackslash}p{(\linewidth - 4\tabcolsep) * \real{0.1300}}
  >{\raggedright\arraybackslash}p{(\linewidth - 4\tabcolsep) * \real{0.4800}}
  >{\raggedright\arraybackslash}p{(\linewidth - 4\tabcolsep) * \real{0.3900}}@{}}
\caption{Result families and evidence boundaries}\label{tab:family-map}\\
\toprule\noalign{}
\begin{minipage}[b]{\linewidth}\raggedright
Family
\end{minipage} & \begin{minipage}[b]{\linewidth}\raggedright
Precisely bounded contribution
\end{minipage} & \begin{minipage}[b]{\linewidth}\raggedright
Evidence boundary
\end{minipage} \\
\midrule\noalign{}
\endhead
\bottomrule\noalign{}
\endlastfoot
QB & Bipartite \(\Delta\leq6\), \(e\geq3n-5\) implies a nonempty quartic
& Main theorem formal \\
TWIN & A pair in every nonempty capped host with \(e\geq3n-2\) and no
singleton left neighbourhood class & Extended actual-host criterion
formal \\
FINITE & Exact small values and larger explicit lower constructions &
Formal exactness through seven; other coverage separated \\
REG5 & Pair-free five-regular graphs of orders 18, 32 and 104 & Examples
formal; minimum-order claim separately qualified \\
SUB & Pair-free gadget substitution for \(d\leq7\) and a vertex-deletion
equivalence & Full preservation, regularity, bipartiteness and
equivalence formal \\
HAAR & Four colours force a pair in every prime-field module; specified
composite template & Every exact-four-colour component decomposes
formally; known portions credited \\
ONE & An arbitrarily large family defeats the supplied-colouring
one-round detector & One-round obstruction and actual full-host
two-round sharpness formal \\
BM & Candidate quantitative one-sided spectral Hamiltonicity adaptation
& Written main argument; missing formal infrastructure explicit \\
\end{longtable}
}

Complete technical proofs and the current verification ledger are in the
\href{https://github.com/alejandrozu/openmath-2026-judging/tree/70c5ec3107658182f22c207d4c50f53f9cb3e602/robert-review-2026-10-07}{companion
repository}. Exact named declarations are listed in Section 10. The
repository also retains the longer written dossiers where reproducing
every technical lemma would obscure the principal statements here.

\section{Formal and graph
conventions}\label{formal-and-graph-conventions}

Formal certificates use the upstream predicate
\texttt{HasTwoEdgeDisjointCyclesSameVertexSet}: closed Lean graph walks
that satisfy \texttt{IsCycle}, equal sets of visited vertices, and
disjoint lists of traversed edges. Robert's finite-set predicate
\texttt{HasPairF} is proved equivalent to it. Connectedness is part of
the cycle predicate; equal degree sequences or disconnected two-factors
are insufficient.

In the substitution section, the skeleton is a loopless labelled
multigraph. Parallel edges remain distinct, and two parallel edges form
a cycle of length two. The substituted graph itself is simple. In Haar
graphs, the two shores are distinct copies of the underlying additive
group, even when their coordinate values coincide. Recolouring changes
edge labels, never host edges or vertex identities.

The frozen public source is Robert Huynh's revision
\texttt{1efc324e155ef0b6a7081c5f695c6c825e7debef}. The previous audit
compiled 106 public authored source files and checked 84 distinct
selected declarations against the standard kernel axiom set. New proofs
in this edition are additive and have their own source hashes and actual
compiler receipts. Compilation does not establish novelty, and imported
published theorems must retain their precise hypotheses. Unresolved
target placeholders and unavailable private projects are excluded from
certified claims.

\section{Quartic extraction below the classical
threshold}\label{quartic-extraction-below-the-classical-threshold}

\textbf{Theorem Q.} Let \(G\) be a finite simple bipartite graph with
\(n\geq3\), maximum degree at most six and \(e(G)\geq3n-5\). Then \(G\)
contains a nonempty four-regular subgraph.

The theorem is
\href{https://github.com/alejandrozu/openmath-2026-judging/blob/70c5ec3107658182f22c207d4c50f53f9cb3e602/robert-review-2026-10-07/lean/Openmath/Proofs/QB5/Statement.lean\#L42}{\texttt{Erdos585.qb5}};
\texttt{qb4}, at \(3n-4\), is a weaker earlier version. The threshold is
independent of a particular labelled graph or finite census. Robert's
proof is a universal Lean proof, including a bounded case analysis
evaluated by the kernel.

The closest located comparator is Alon, Friedland and Kalai {[}1, Remark
3.6{]}. For bipartite graphs, their divisibility argument yields a
nonempty four-divisible subgraph once \(e>3(n-1)\). With maximum degree
at most seven, nonzero degrees in that subgraph must be four, so the
integer threshold is \(3n-2\). Theorem Q lowers it by three edges under
the stronger maximum-degree-six assumption. One must not transfer
results whose proof selects seven matching classes directly into the
degree-six setting.

\textbf{Proof of Theorem Q.} Here is the proof chain of the four
original QB5 papers, including its computer-assisted boundary. For a
vertex set \(S\), write \(g(S)=6|S|-2e(S)\). Choose a counterexample
with first \(n\), then its edge count minimal. It has \(e=3n-5\), hence
\(g(V)=10\). For every proper \(S\) with \(|S|\ge3\), minimality gives
\(g(S)\ge12\); deleting a vertex therefore gives minimum degree at least
four. The cases \(3\le n\le9\) are excluded by the bipartite edge bound,
and the only ten-vertex possibility is \(K_{5,5}\), which contains a
quartic.

If the shore sizes differ by \(j\ge0\), and the smaller shore has
deficiency \(d=\sum(6-\deg v)\), edge counting gives \(10=6j+2d\). Thus
either the shores are balanced with deficiency five each (E5), or their
sizes differ by one with deficiencies two and eight (C2). Part I's C2
hub-and-petals lemma produces a vertex on the larger shore whose
deletion has a four-factor. Its hypotheses are the preceding sparsity
inequalities, not an assumed factor; the analytic proof is implemented
by the C2Hub/C2Petals/C2TwoBlock modules.

It remains to exclude E5. The bipartite four-factor flow criterion
supplies a violating cut. Its slack, together with \(g(S)\ge12\), forces
a decomposition

\[
X=A\sqcup C,\qquad Y=B\sqcup D,\qquad
|A|=|C|+2,\quad |D|=|B|+2.
\]

Vertices of \(C\) are saturated inside \(A\), and vertices of \(B\)
inside \(D\). Exactly seven edges join \(A\) to \(D\). Both blocks have
\(g=12\). For a selected four-set \(M\) of those edges, let \(m(T)\)
count its ends in \(T\subseteq A\), and define

\[
\operatorname{dem}_X(T)=4|T|-\sum_{c\in C}\min(4,e(c,T)),\qquad
I_X(M)=\bigcap_{m(T)<\operatorname{dem}_X(T)}T.
\]

The family in this intersection is nonempty: \(A\) has demand eight but
supply four. The flow criterion says that deleting \(p\in A,q\in D\)
leaves a four-factor precisely when some \(M\) satisfies
\(p\in I_X(M)\), \(q\in I_Y(M)\). Indeed these memberships imply no
selected edge meets \(p,q\), and all demands on the remaining vertices
are met. Conversely the four-factor uses exactly four cut edges, and its
degree sums give those demand inequalities.

Part IV proves KL1: if \(M\) meets every in-block degree-three vertex,
then \(I_X(M)\ne\varnothing\). Its mechanism is uncrossing. Put
\(\delta_a=6-\deg_X a\), \(w(a)=\delta_a-m(a)\); then \(w(A)=8\),
\(0\le w(a)\le2\). The complement of an under-supplied set is overloaded
when

\[
\theta(T)=m(T)+4-4|T|+\sum_{c\in C}(e(c,T)-2)_+\ge1.
\]

If \(I_X(M)\) were empty, maximal overloaded sets would cover \(A\). The
induced deficiency identity gives supermodularity with a nonnegative
crossing-edge bonus. A minimal cover by maximal overloaded sets must
therefore contain a common hub \(K\), whose enlargement \(K^+\) has
\(\theta=2\). The remaining petals \(P_i\) are disjoint; put
\(Q_i=P_i\cap A\), \(E_i=e(K^+,P_i)\). In the equal-shore hub case,
\(m(K)=4\), there are no singleton members and every petal satisfies
\(w(Q_i)<E_i\). Summing contradicts
\(8-w(K)=\sum w(Q_i)<\sum E_i\le8-w(K)\).

In the other case \(w(K)=0\), \(\sum E_i\le m(K)+6\), and

\[
w(Q_i)\le\tfrac23\bigl(E_i+m(Q_i)\bigr)
+\tfrac13\mathbf1_{\{|P_i|\ge3,\ |Q_i|-|P_i\cap C|=2\}}.
\]

Without a singleton member, weight eight requires at least four
exceptional petals, but these require at least 27 boundary edges against
a budget at most ten. With a singleton, \(m(K)=3\), no exception is
possible, and equality forces three singleton petals. A remaining vertex
of \(C\) would then have degree six but only four possible neighbours.
Parts III--IV prove the precise maximal-family/hub lemmas used here;
their complete analytic proofs remain in the mathematical supplement.

A chosen \(M\) need not meet every degree-three vertex. The final
\textbf{seven-edge covering lemma} supplies a four-set missing at most
one such vertex on each side. It also excludes every rigid set at a
missed vertex: a set of selected ends whose marked members have cut
degree three, all its incident cut edges lie in \(M\), and their number
is three or four. Part III's rigidity lemma makes such a set necessary
if the missed vertex is overloaded. Thus each side has a vertex in
\(I(M)\), either by KL1 or by rigidity exclusion, and the factor just
constructed contradicts minimality.

This last covering lemma is a finite certificate, not a claimed
closed-form ordinary argument. It encodes the endpoint partitions of
seven labelled edges, admissible degree-three marks and the 35
four-subsets. The encoding-to-graph proof and kernel evaluations
establish the exact two-sided conclusion \cite{ref11}. Neither the earlier
large graph census nor the ``\(n\ge40\)'' computational observation is
needed. This completes the certificate-assisted proof of Q \cite{ref11}.
\textbf{Corollary Q1.} If \(B\) is a bipartite six-regular graph and
\(v\) is any vertex, then \(B-v\) contains a nonempty four-regular
subgraph. Indeed, if \(B\) has \(N\) vertices, deletion gives \(n=N-1\)
and \(e=3N-6=3n-3\geq3n-5\), while preserving bipartiteness and the
degree cap. This corollary does not assert that \(B-v\) contains a pair.

The candidate original contribution is the precise improved threshold
and its proof mechanism. This is not a new proof of the existence of
regular subgraphs in general, nor a Hamilton decomposition theorem. The
sharpness of \(3n-5\) and the possible extensions to \(3n-6\) or
\(3n-7\) remain separate questions.

\section{Repeated neighbourhoods force an actual
pair}\label{repeated-neighbourhoods-force-an-actual-pair}

\textbf{Theorem T.} Let \(G\) be a nonempty finite simple graph on
disjoint shores \(A\sqcup R\), with edges only between shores, maximum
degree at most six and \(e(G)\geq3|V(G)|-2\). Suppose every \(a\in A\)
has a distinct vertex \(b\in A\) with exactly the same neighbourhood in
\(R\). Then \(G\) contains a pair.

The frozen source proves the exact-count case. The added endpoint
\href{https://github.com/alejandrozu/openmath-2026-judging/blob/70c5ec3107658182f22c207d4c50f53f9cb3e602/robert-review-2026-10-07/lean/TwinDensityExtension.lean\#L132}{\texttt{OpenMathReview.TwinDensityExtension.hasPair\_of\_no\_singleton\_twins\_dense}}
proves the full inequality range. Hypotheses are read directly from
\(G\); no quotient cycle or successful pair is assumed. Nonemptiness is
essential because the empty graph satisfies \(e+2\geq3n\) without
containing a pair.

\textbf{Proof.} Put \(a=|A|\), \(r=|R|\). The two nonnegative shore
deficits \(6a-e\) and \(6r-e\) sum to at most four and differ by
\(6(a-r)\), so \(a=r\) and their common value \(d\) is at most two. A
nonzero deficit from a repeated-neighbourhood class has multiplicity at
least two, ruling out \(d=1\). A class of size at least four must have
six common neighbours, or its deficit would exceed two; it gives
\(K_{4,4}\).

Otherwise let \(p,q\) count the classes of sizes two and three and let
\(g=p+q\). The quotient incidence graph has left degrees at most six and
right degrees at most three. Write \(m\) for its number of incidences.
If \(d=0\), then \(m=6g\) and \(r=2p+3q\), so \(m\geq3g+r\). If \(d=2\),
exactly one pair class loses one common neighbour; hence \(m=6g-1\) and
\(p\geq1\), giving the same inequality.

Apply the classical zero-sum selection {[}1, Theorem 2.1{]} in
\((\mathbb Z/4\mathbb Z)^g\oplus(\mathbb Z/2\mathbb Z)^{r-1}\). The
bound \(m\geq3g+r\) gives a nonempty selected incidence set whose left
degrees are zero or four. The omitted right parity follows from the even
selected total, so right degrees are zero or two. Take an active
connected component and suppress the degree-two right vertices. The
result is a connected four-regular loopless multigraph with individually
labelled edges. Its subdivided double embeds in \(G\) using two actual
twins in each active class. The known Euler-tour lift gives two
complementary Hamilton cycles on that subgraph \cite{ref8}; their actual
supports and edges map faithfully into \(G\). ∎

The selected subgraph is precisely a subdivided double, and Eppstein's
prior is stronger: every Hamilton cycle there has a Hamiltonian
complement. The selection is classical Olson/AFK zero-sum theory. The
possible contribution is the exact actual-host criterion and its
formalization, rather than a new counting principle or Eulerian lift.
Its priority remains unresolved and it may be a routine corollary of
those ingredients. The at-most-two-singleton extension is proved below.

The following is an ordinary source-derived corollary within the
existing TWIN family. It is not represented as a separately selected
Lean endpoint or a new family credit.

\textbf{Corollary.} Every nonempty finite simple bipartite six-regular
graph whose one shore has at most two singleton maximal
equal-neighbourhood classes contains a pair.

\textbf{Proof.} Six-regularity balances the two shores. A twin class of
size at least four immediately supplies \(K_{4,4}\), using four of its
six common neighbours. Otherwise let \(p,q,s\) count maximal classes of
sizes two, three and one. Collapse the nonsingleton classes and omit the
singleton vertices while retaining all right vertices. The quotient has
\(g=p+q\) group vertices, \(r=2p+3q+s\) right vertices and \(m=6(p+q)\)
incidences. Its right degrees are at most three. The parity-dependent
mixed-modulus selection criterion requires \(m>3g+r-1\), which holds
when \(p\ge s\). It selects group degrees zero or four and right degrees
zero or two. Suppressing the active right vertices produces a labelled
loopless four-regular multigraph; an Euler circuit and two distinct
actual twins in each active group lift to two edge-disjoint connected
cycles with identical support.

For \(s=0\), \(p\ge s\) is automatic. For \(s=1\), a neighbour of the
singleton has five other neighbours, consisting of entire twin classes.
With class sizes two or three, five can only be partitioned as \(2+3\),
so \(p\ge1\). For \(s=2\), suppose \(p\le1\). A common neighbour of the
two singletons would need four remaining neighbours from two distinct
pair classes, impossible. The two six-neighbour sets are therefore
disjoint. Each of those twelve right vertices needs a pair class among
its remaining five neighbours. If \(p=0\) this is impossible; if
\(p=1\), the unique pair class would have twelve neighbours despite
degree six. Thus \(p\ge2\). In all cases \(s\le2\) implies \(p\ge s\),
and the selection-and-lifting argument applies. \(\square\)

There is no claim for three singleton classes, no assumption that every
six-regular graph has such twins, and no deduction of B6. The ordinary
proof and its review are in
\href{https://github.com/roberthuynh/erdos-585/blob/1efc324e155ef0b6a7081c5f695c6c825e7debef/findings/supplement/papers/twin-core/TWIN-ZEROSUM.md}{\texttt{TWIN-ZEROSUM.md},
Section 5}. The no-singleton formal density theorem and this ordinary
restricted six-regular extension should be distinguished.

\section{Finite extremal profiles and lower
constructions}\label{finite-extremal-profiles-and-lower-constructions}

The source reports the profile

\[
(f(1),\ldots,f(12))=(0,1,3,6,9,12,16,19,23,27,31,36).
\]

The evidence is not uniform across this row. Exact values through seven
have formal statements; the finite procedures through ten were
independently replayed in the prior audit. The displayed exact upper
exclusions at eleven and twelve remain source-reported computational
coverage, while the lower bounds \(f(11)\geq31\) and \(f(12)\geq36\)
have independently compiled construction proofs. We do not promote a
lower witness into an exact maximum.

\textbf{Basic counting.} If a pair has support of size \(s\), its two
cycles use \(2s\) distinct edges, hence \(2s\leq\binom{s}{2}\) and
\(s\geq5\). Therefore \(f(n)=\binom n2\) for \(n\leq4\), and \(K_5\)
minus an edge is pair-free. Adding a new vertex with three earlier
neighbours preserves pair-freeness: any pair containing that vertex
would require four distinct incident edges. This supplies useful linear
extensions without claiming an asymptotic improvement.

The source's larger formal lower bounds should remain visible rather
than being reduced to the phrase ``linear constructions'':

Table~\ref{tab:lower-bound-map} records the explicit lower-bound constructions and their parameter domains.

{
\begin{longtable}[]{@{}
  >{\raggedright\arraybackslash}p{(\linewidth - 2\tabcolsep) * \real{0.4900}}
  >{\raggedright\arraybackslash}p{(\linewidth - 2\tabcolsep) * \real{0.5100}}@{}}
\caption{Explicit lower-bound constructions and parameter domains}\label{tab:lower-bound-map}\\
\toprule\noalign{}
\begin{minipage}[b]{\linewidth}\raggedright
Bound
\end{minipage} & \begin{minipage}[b]{\linewidth}\raggedright
Parameter domain and certificate
\end{minipage} \\
\midrule\noalign{}
\endhead
\bottomrule\noalign{}
\endlastfoot
\(f(n)\ge4n-13\) & \(n\ge10\); the two-hub sparse-rim family \\
\(f(11)\ge31,\ f(12)\ge36\) & Two-hub witness and two adjacent hubs
joined to the Petersen graph \\
\(f(2k^2+2)\ge8k^2-3k+1\) & \(k\in\mathbb{N}\); complete-base matching
subdivision \\
\(f(m^2+3m+2)\ge5m^2\) & \(m\in\mathbb{N}\); Latin incidence with two
exceptional hub vertices \\
\(f(n)\ge5n-15\lfloor\sqrt{n}\rfloor-10\) & \(n\ge16\); selected Latin
triples and padding \\
\(f(n)\ge3n-6,\ 3n-5,\ 3n-4\) & Respectively \(n\ge5,7,9\); the smaller
hub/rim witnesses \\
\(f(n+1)\ge f(n)+3\) & \(n\ge3\); adjoining a vertex with three
neighbours \\
\end{longtable}
}

Here the inequalities are about the actual maximum over pair-free simple
graphs. In the Lean source subtraction is natural subtraction; the
stated parameter restrictions prevent a misleading negative
interpretation of the uniform bounds. The formal statements are in
\texttt{Openmath/Proofs/Final.lean}, especially lines 102--126. None
changes the established asymptotic lower order \(\Omega(n\log\log n)\).

For \textbf{matching subdivision}, let \(H\) have \(v\) vertices and
\(e\) edges, and let \(M\subseteq H\) be a matching with \(m\) edges.
Keep the edges of \(M\), subdivide each remaining edge once, and adjoin
two adjacent hubs joined to every other vertex. The actual construction
has

\[N=v+e-m+2,\qquad E=2v+4e-3m+1.\]

\textbf{Why matching subdivision is pair-free.} Suppose its union
contained a nonempty four-regular subgraph. An ordinary base vertex has
at most one retained matching neighbour and two hub neighbours, so some
subdivision vertex must occur. Every subdivision vertex has exactly four
neighbours, forcing both hubs and its two base endpoints into the
subgraph. Let \(t\) count selected subdivision vertices, \(v\) selected
base vertices, and \(h\in\{0,1\}\) indicate use of the hub edge. The two
hubs' degree sum gives \(8-2t-2h\) hub-to-base edges. Counting the base
vertices' degrees then gives \(4v\le(8-2t-2h)+2t+v\), because the
retained matching contributes at most \(v\) edge ends. Hence \(v\le2\),
while a subdivision vertex forces \(v\ge2\). With only two base
vertices, simplicity permits at most one subdivision vertex, and its
base pair cannot also be a retained matching edge. Each base vertex
would consequently have degree at most three, a contradiction. A pair
would have a four-regular union, so none exists.

The pair-exclusion theorem uses the matching degree bound; the displayed
counting identities additionally require \(M\subseteq H\). These
hypotheses must not be conflated. Taking \(H=K_{2k}\) and a perfect
matching gives the bound in the table. At \(k=2\) it gives the
ten-vertex, 27-edge lower witness. The proof works with genuine cycles
and actual subdivision vertices, not merely a claimed four-factor.

For \textbf{Latin incidence}, use triples from the addition table of a
cyclic group of order \(m>0\). There are \(m^2\) line vertices and
\(3m\) point vertices. Each line is incident with exactly three points,
and two distinct ordinary points lie on at most one common line. Put two
further hub vertices on the point shore and join each to every line
vertex. There are no edges between the hubs and no added hub-to-point
edges: this differs from the adjacent-universal-hub matching
construction. Each line has degree five, giving \(N=m^2+3m+2\) and
\(E=5m^2\). The incidence exclusion criterion uses the
three-ordinary-neighbour cap and uniqueness of the line through two
ordinary points. It is not valid for an arbitrary incidence relation
lacking those conditions. \textbf{Why Latin incidence is pair-free.} A
four-regular subgraph must contain at least one hub, since a line has
only three ordinary neighbours. With one hub, its degree forces exactly
four selected lines; regularity balances the bipartition, leaving three
ordinary point vertices. Every selected line would contain all three
points, contrary to uniqueness of the line through two ordinary points.
With two hubs, let \(\ell\) count selected lines. There are \(\ell-2\)
ordinary point vertices and exactly eight hub-to-line edges. Each line
uses one or two hub edges, so \(4\le\ell\le8\). Exactly \(8-\ell\) lines
use two ordinary points and \(2\ell-8\) use three. They therefore use
\(5\ell-16\) distinct ordinary-point pairs. Uniqueness allows at most
\(\binom{\ell-2}{2}\) pairs, which is strictly smaller for each integer
\(4\le\ell\le8\). This again excludes a quartic, and hence a pair.

Selecting a subset of the triples and adding harmless low-degree
vertices gives the uniform square-root-error bound. The case \(m=0\) of
the lower inequality is handled separately in the source.

The exact finite profile through twelve remains

\[ (0,1,3,6,9,12,16,19,23,27,31,36). \]

Formal exactness through seven, fresh computational coverage through
ten, and source-reported upper exclusions at eleven and twelve are
different evidence classes. The two formal lower witnesses at eleven and
twelve establish attainment, not the exhaustive upper inequalities.
Likewise, an edge census or a regular-graph census must retain its
generator completeness, replication and stopping limitations. The
publication contribution is a reproducible finite record and any newly
identified construction mechanism. The trivial small values and routine
extension argument are background; a complete minimality or extremality
claim requires both attainment and exhaustive upper coverage. The linear
construction families do not surpass the known \(\Omega(n\log\log n)\)
lower order. A reusable certificate implementation can still be valuable
even where the underlying numerical value is already known.

\section{Five-regular pair-free
examples}\label{five-regular-pair-free-examples}

\textbf{Theorem W.} There exist simple five-regular pair-free graphs on
18 and 32 vertices, and a simple bipartite five-regular pair-free graph
on 104 vertices.

The three five-regular examples certify existence at orders 18, 32 and
104; they do not constitute three asymptotic improvements. Their
mechanisms are worth stating explicitly because regularity alone hides
the reason a pair is excluded. Each vertex on the common support of a
pair has four incident edges in the union. Moreover, if that support
meets both sides of a cut, each connected cycle crosses the cut an even
positive number of times. Two edge-disjoint cycles therefore use at
least four distinct cut edges. A cut of size at most three confines
every pair to one side.

\textbf{The eighteen-vertex witness.} Start with a nine-vertex block.
Vertices \(0,1,2\) form a triangle and are each joined to \(3,4,5\).
Vertices \(6,7,8\) form another triangle; add
\(6\!-\!3,6\!-\!4,6\!-\!5,7\!-\!3,7\!-\!4,8\!-\!5\). There are 21 edges.
Vertices 7 and 8 have degrees four and three, respectively, and every
other vertex has degree five. In the displayed order each vertex has at
most three earlier neighbours, so the block is three-degenerate and
cannot contain a pair. Take two copies, with the second labelled
\(9,\ldots,17\), and add \(8\!-\!17,8\!-\!16,7\!-\!17\). These links
fill precisely the two copies' degree deficiencies. The resulting graph
has 18 vertices and 45 edges and is five-regular. Its three-edge joining
cut confines any alleged pair to a pair-free block. This example can
have quartic subgraphs crossing the cut in two edges; quartic extraction
still does not supply a Hamilton decomposition. The source gives an
explicit 17-vertex quartic support illustrating that distinction.

\textbf{The thirty-two-vertex witness.} A block consists of the double
wheel on a five-cycle with its hub edge removed, together with a
degree-three vertex. Two such eight-vertex blocks form a half, joined by
three edges; two halves are joined by two edges. The exact labelled
joins are given in the source edge list. Regularity is checked on the
actual graph, which has 32 vertices and 80 edges. The two-edge cut first
confines a pair to one half, and the three-edge cut confines it to one
block. A degree-three vertex cannot belong to the pair's support; the
remaining seven-vertex double wheel is itself pair-free. To see this, a
quartic support must use both hubs. Every selected rim vertex must then
use both of its rim neighbours, forcing the entire five-cycle into the
support; each hub would have degree five instead of four. Thus the
double wheel has no quartic support. Thus the exclusion is a hierarchy
of small-cut and local-support certificates, not an exhaustive search of
all five-regular graphs.

\textbf{The bipartite 104-vertex witness.} The thirteen-vertex block has
six inner and seven outer vertices. It is built from \(K_{3,3}\) by
adjoining seven vertices with three earlier neighbours. Its
three-degeneracy excludes a quartic subgraph and hence a pair. Every
inner vertex has degree five; the outer deficiencies are \(1,2,2\) at
three ports. Eight copies are linked by 20 edges according to the
specified five-regular bipartite multigraph skeleton, filling all
deficiencies and producing 104 vertices and 260 edges. The hierarchy of
two- and three-edge cuts confines a pair to a block. The published
formal endpoint verifies the actual simple graph, bipartition,
regularity and pair exclusion directly by this cut route. The
construction also fits the gadget interpretation, but its original proof
is not silently replaced by an assumed projected-cycle argument. A
second, nonisomorphic port assignment is an author-supplied
computational example in the same family; the original formal endpoint
covers the first assignment.

The separate claim that 18 is the least possible order depends on the
complete exclusion census through 16 vertices. One complete method and a
partial independent second method are documented; the latter is not a
second full census. We retain the claimed minimum as source-reported
computational evidence, separately from the three independently compiled
existence statements. Generic nonbipartite five-regular pair-free
existence has known predecessors, so witness identities, minimality and
bipartite priority require their own comparisons. Generic non-bipartite
five-regular pair-free existence is classical: Read and Wilson {[}9,
p.155{]} record five-regular graphs without a quartic subgraph, which
therefore have no pair. This does not identify the eighteen-vertex
witness, prove its minimum order, or settle bipartite pair-freeness. A
five-regular bipartite graph always has a quartic factor by removing a
perfect matching, so the bipartite example needs actual cycle-pair
exclusion. Exact witness and construction priority remain separate
questions.

\section{Substitution and prescribed-vertex
deletion}\label{substitution-and-prescribed-vertex-deletion}

\textbf{Definition.} A \(d\)-gadget is a finite simple bipartite graph
\(\Gamma\) with shores \(I,O\), every inner vertex of degree \(d\),
every outer vertex of degree at most \(d\), and \(|O|=|I|+1\). The total
outer deficiency is \(d\). A valid substitution into a loopless
\(d\)-regular multigraph \(L\) replaces every skeleton vertex by a copy
of \(\Gamma\) and every labelled skeleton edge by one link between outer
ports. Each port receives exactly its deficiency in links, and the final
graph is simple.

\textbf{Theorem S (formal preservation theorem).} If \(d\leq7\), both
\(\Gamma\) and \(L\) are pair-free, and \(G\) is a valid substitution,
then \(G\) is pair-free and \(d\)-regular. If \(L\) is bipartite, so is
\(G\).

\textbf{Proof.} Suppose a pair has common support \(S\). In one gadget
copy, put \(S_I=S\cap I\) and \(S_O=S\cap O\), and let \(x\) count its
used link edges in the union of the two cycles. Every selected vertex
has union degree four. Counting internal edge ends on the two shores
gives

\[
x=4(|S_O|-|S_I|).
\]

There are at most \(d\leq7\) links at that copy, so \(x\) is zero or
four. If it is zero, connectedness confines both cycles to the same
pair-free gadget. Otherwise both cycles meet the copy and its
complement, and each uses an even positive number of cut edges. Four in
total therefore means exactly two per cycle. Contract the gadget copies
along either cycle: the retained labelled link edges form a connected
two-regular loopless multigraph on the same set of visited copies. They
are a cycle, including the possible length-two case. The two projected
cycles are edge-disjoint because the link correspondence is a bijection.
This contradicts pair-freeness of \(L\). Regularity follows from filling
deficiencies; the bipartition follows by reversing gadget shores on one
side of the skeleton. ∎

The degree bound is the single-passage threshold in this proof. For
eight available links, the counting identity also allows \(x=8\), so a
two-passage pattern is no longer excluded by this argument. This is a
boundary of the method, not a counterexample to every possible
substitution theorem at larger degree.

\textbf{Corollary S1 (formal prescribed-vertex equivalence).} Let B6
assert that every nonempty finite simple bipartite six-regular graph
contains a pair. Then B6 is equivalent to the statement that, for every
such graph \(B\) and every vertex \(o\), the graph \(B-o\) contains a
pair. For the nontrivial direction, if \(B-o\) is pair-free, it is a
six-gadget whose six distinct neighbours of \(o\) each have deficiency
one. Substitute it into the supplied pair-free bipartite six-regular
multigraph skeleton. Distinct unit-deficiency ports make the resulting
substitution simple; Theorem S gives a bipartite six-regular pair-free
graph, contradicting B6. The opposite implication is immediate.

The new single endpoint
\href{https://github.com/alejandrozu/openmath-2026-judging/blob/70c5ec3107658182f22c207d4c50f53f9cb3e602/robert-review-2026-10-07/lean/R1.lean\#L13}{\texttt{RobertPublishable.SUB.theorem\_R1}}
proves pair-freeness, actual degree \(d\) and bipartiteness from the
faithful gadget/link/port data. It derives connected projected cycles
and counts parallel edges separately. Its preservation proof even
permits heterogeneous pair-free copy fibres with cut size at most seven,
and a locally finite infinite skeleton; cycles remain finite closed
walks. The endpoint
\href{https://github.com/alejandrozu/openmath-2026-judging/blob/70c5ec3107658182f22c207d4c50f53f9cb3e602/robert-review-2026-10-07/lean/PrescribedVertex.lean\#L89}{\texttt{RobertPublishable.SUB.prescribed\_vertex\_equivalence}}
also constructs and verifies the explicit sixteen-vertex,
forty-eight-edge pair-free bipartite six-regular multigraph skeleton and
the actual six distinct unit ports of the deleted gadget. Both sides
quantify over finite nonempty hosts. Neither assertion proves B6.

\section{Haar graphs and a four-colour
threshold}\label{haar-graphs-and-a-four-colour-threshold}

Let \(V\) be an additive group and \(S\subset V\) a finite palette. The
Haar graph \(H(V,S)\) has shores \(V\times\{0\}\) and \(V\times\{1\}\),
with \(x_0\) adjacent to \((x+s)_1\) for each \(s\in S\). Every colour
is a translation matching. These graphs are not automatically ordinary
Cayley graphs on an abelian group.

\textbf{Theorem H.} Let \(p\) be prime and let \(V\) be any
\(\mathbb F_p\)-module. If \(|S|\geq4\), then \(H(V,S)\) contains a
pair. Conversely, a pair requires at least four colours. The forward
theorem is formal without finiteness of \(V\); the new
\href{https://github.com/alejandrozu/openmath-2026-judging/blob/70c5ec3107658182f22c207d4c50f53f9cb3e602/robert-review-2026-10-07/lean/HaarInfinite.lean\#L110}{\texttt{OpenMathReview.HaarInfinite.hasPair\_iff}}
completes the converse under the same unrestricted ambient-module
hypotheses. This is a finite cycle-pair assertion, not an infinite
Hamilton cycle.

\textbf{Hamilton decomposition corollary.} If \(|S|=4\), every connected
component of \(H(V,S)\) is finite and decomposes into two Hamilton
cycles. Given a pair, its union supplies four distinct neighbours at
every support vertex, exhausting the four host neighbours. Its support
is therefore closed under all host edges. Connectedness makes that
support the whole component. Translations and the shore-reversing map
\((x,b)\mapsto(-x,1-b)\) are actual graph automorphisms, so a pair can
be moved to any vertex. The new endpoint
\href{https://github.com/alejandrozu/openmath-2026-judging/blob/70c5ec3107658182f22c207d4c50f53f9cb3e602/robert-review-2026-10-07/lean/HaarComponents.lean\#L73}{\texttt{RobertPublishable.HaarComponent.every\_connected\_component}}
formalizes both the spanning cycles and the exhaustion of component
edges. This is a strengthened formulation of the four-colour result, not
an additional independent discovery. For more than four colours, the
pair lies in a selected four-colour subgraph; the conclusion does not
decompose the larger-palette host.

\textbf{Proof of the four-colour construction.} For a finite coordinate
set, a matching permutation \(P\) sends \(x_L\) to \(P(x)_R\). Two
edge-disjoint matchings \(P,Q\) form one Hamilton cycle exactly when
\(Q^{-1}P\) has one orbit. In \(\mathbb F_p^2\), a skew shift with
nonzero horizontal step has one orbit if its total vertical increment
over the \(p\) horizontal positions is nonzero: \(p\) steps return
horizontally and translate vertically by that total.

\textbf{Rank two.} Normalize the colours to
\(0,(1,0),(0,1),(\alpha,\beta+1)\), with \(\alpha\ne0\) and the fourth
colour different from \((1,0)\). The four matchings are

\[
\begin{aligned}
R_0(x,y)&=(x+1,y),&R_1(x,y)&=(x,y+\mathbf1_{x=0}),\\
B_0(x,y)&=(x+\alpha,y+\beta+1),&
B_1(x,y)&=(x,y+\mathbf1_{x\ne0}).
\end{aligned}
\]

They partition the host: the two vertical choices use complementary
zero/unit colours, while the other translations are distinct. The red
successor has vertical increment \(-\mathbf1_{x+1=0}\), totalling
\(-1\). The blue successor has increment
\(\beta+\mathbf1_{x+\alpha=0}\), totalling \(p\beta+1=1\). Both are full
orbits.

If the fourth colour is \((0,v)\), \(v\ne0,1\), use \(F(x,y)=(-x-y,y)\),
translating the right shore by \((1,0)\). This maps the palette
\(0,(1,0),(0,1),(1-v,v)\) bijectively to the desired palette. Since
\(1-v\ne0\), the preceding construction applies \cite{ref11}.

\textbf{Rank three.} For colours \(0,e_a,e_b,e_c\), use \(A(x)=x+e_a\),
\(C(x)=x+e_c\), and

\[
Z(a,b,c)=(a,b+\mathbf1_{a=c},c),\qquad
B(a,b,c)=(a,b+\mathbf1_{a\ne c},c).
\]

The red successor \(Z^{-1}A\) has one \(p^2\)-point orbit in each
fixed-\(c\) level: its vertical period increment is \(-1\). The blue
successor \(B^{-1}C\) similarly has one orbit per fixed-\(a\) level,
with increment \(-(p-1)=1\). Mark \(m_u=(u,0,-u)\). Exchange \(A\) and
\(C\) at these marks. This preserves both permutations because
\(A(m_u)=C(m_{u+1})\). The new red successor joins the cut level paths
in order \(u\mapsto u-1\), and blue joins them in order
\(u\mapsto u+1\). Each now has one orbit on all \(p^3\) coordinates.
Their matching unions partition the host, giving two Hamilton cycles on
the same \(2p^3\) vertices \cite{ref11}.

Rank one uses two distinct translation pairs, each with a nonzero
prime-field step. Selecting four colours, subtracting an anchor and
using an injective rank-one/two/three coordinate chart exhausts the
possibilities. Add the anchor only on the right shore. These maps
preserve actual edges and supports; they work in an infinite ambient
module because the selected chart is finite. The prime-line case is
covered by the dihedral decomposition theorem
\href{https://arxiv.org/html/1810.07866v1}{the prime-dihedral preprint,
Theorem 3}. Degree-four abelian Cayley decomposition \cite{ref6} covers
characteristic two and centrally symmetric odd-prime palettes: for
\(S=\{c\pm u,c\pm v\}\), shifting the right shore by \(-c\) identifies
the four-colour component with an abelian Cayley graph on the subgroup
generated by \((\pm u,1),(\pm v,1)\) in \(V\times\mathbb Z_2\). Neither
comparison settles every nonsymmetric odd-prime rank-two or rank-three
palette.

The recent generalized-dihedral preprint
\href{https://arxiv.org/html/1810.13311v3}{the current
generalized-dihedral preprint} proves one Hamilton cycle; the
cyclic-Haar result {[}10, Proposition 5.1{]} also proves one cycle under
its factorization hypotheses. Neither is a decomposition theorem. Thus
the candidate new portion is the residual complementary-cycle
construction and the specified composite template, subject to exact
priority review.

\textbf{Composite template and certificate.} The separate cyclic
construction has ambient group \(\mathbb Z/n^2\mathbb Z\), palette
\(\{0,n,2n,1\}\), \(n\ge3\). Write \(r(x)=x\bmod n\), and mark the
integer representatives \(0,\ldots,n-1\). Let \(\sigma_\pm\) shift these
marks by \(\pm1\bmod n\) and fix every other coordinate. Put

\[
R_1(x)=\begin{cases}x&r(x)=0,\\x+n&r(x)\ne0,\end{cases}
\quad
S(x)=\begin{cases}x+n&r(x)=0,\\x&r(x)\ne0.\end{cases}
\]

The four actual matchings are \(R_0=S\sigma_+\), \(R_1\),
\(B_0=(x\mapsto x+1)\sigma_-\), \(B_1=(x\mapsto x+2n)\). Their values
use four distinct permitted shifts. The red base successor translates
each residue class by \(n\) or \(-n\), with exactly one mark per class.
The cyclic mark splice joins all \(n\) classes.

For odd \(n\), blue agrees off the marks with translation by \(d=1-2n\),
a unit modulo \(n^2\) with inverse \(1+2n\). Put \(h=(n+1)/2\bmod n\).
In this translation's cyclic order the marks are
\(\operatorname{mark}(hj)\), at integer positions
\(nj+\operatorname{val}(hj)\). After rewiring, the path departing mark
\(h(j+2)\) reaches mark \(h(j+1)\) along the intervening unmarked
interval. Hence all marks lie in one blue orbit; every unmarked point
reaches a mark by following the old unit translation. This proves blue
connectedness without assuming it \cite{ref11}.

For even \(n\ge4\), the old blue successor has two orbits. Set
\(r=n/2-1\) and

\[
a=r_L,\ b=(r+1)_R,\ c=(r+1)_L,\quad
d=(r+1+n)_R,\ e=(r+n)_L,\ f=(r+2n)_R.
\]

Exchange the red edges \(ab,cd,ef\) with the blue edges \(bc,de,fa\).
Each of these six distinct vertices loses and gains one edge in each
factor, so degrees and disjointness are preserved. \textbf{Connectedness
requires more:} the supplied interval-path certificate proves a
surviving red \(b\)-to-\(d\) path, surviving blue \(a\)-to-\(c\) and
\(b\)-to-\(f\) paths, the exact two-orbit blue partition, and global
connectedness of both exchanged factors. These are proved obligations,
not premises inferred from degree two \cite{ref11}. The two connected
degree-two factors give actual spanning cycles. This paragraph explains
that formal certificate; it does not claim a new independent ordinary
proof of its interval lemmas or an arbitrary composite-group
classification.

\section{A recolouring barrier and the quantitative expander
route}\label{a-recolouring-barrier-and-the-quantitative-expander-route}

\subsection{One component-union
round}\label{one-component-union-round}

Let \(q=3^k\) with \(k\geq2\), let \(F=\mathbb F_q\), and take the Haar
graph on \(F^2\) with palette \(P=\{(a,a^2):a\in F\}\). Colour each
translation matching by its parameter \(a\). The graph is simple,
connected, bipartite and four-cycle-free, has \(N=2q^2\) vertices and
degree \(q\).

An allowed \emph{round} chooses two colours and swaps them on any union
of connected components of their bichromatic graph. It may change
arbitrarily many edges or components. The detector examines all
bichromatic cycles exposed from the supplied canonical colouring by zero
or one such round, and may compare cycles from different resulting
states.

\textbf{Theorem R.} Every exposed cycle has induced host degree at most
three on its support. It consequently has no edge-disjoint host cycle on
that same support. For every fixed \(C>0\) and \(A\geq0\), this
obstruction occurs at arbitrarily large orders with \(q>C(\log N)^A\).

\textbf{Proof.} The parabola is a Sidon set, and in characteristic three
this is equivalent to a 2-cap: every subset of at most four points is
affinely independent {[}7, Theorem 3.2{]}. Both this equivalence and the
parabola construction predate the present work. For a nonempty set \(T\)
of at most three colours, choose \(a_0\in T\) and put
\(U=\operatorname{span}_{\mathbb F_3}\{v_a-v_{a_0}:a\in T\}\). A
component has shores \(X+U\) and \(X+v_{a_0}+U\): two-edge walks
generate the differences, and selected edges stay in those cosets. A
fourth palette point in this affine coset would contradict the 2-cap
property. The component is therefore induced in the whole host and has
degree \(|T|\leq3\). After one allowed round, every exposed bichromatic
cycle uses at most three original colours. Its support lies in such an
induced component; a pair would require four distinct incident host
edges at each support vertex, a contradiction. Finally \(q=\sqrt{N/2}\)
exceeds every fixed polylogarithm along this family. ∎

The contribution is the actual recolouring-detector obstruction, not a
new Sidon construction, an avoiding graph of polynomial degree, or a
claim about unrestricted Kempe sequences. The full graph already
contains pairs by Theorem H. Standard Kempe-equivalence results concern
reachability under unrestricted sequences and do not decide this
fixed-radius question \href{https://arxiv.org/abs/2105.01363}{the cited
preprint}.

\textbf{General Sidon-palette theorem.} Let \(V\) be any
\(\mathbb F_3\)-module and \(S\subset V\) a finite Sidon palette. Under
its canonical translation colouring, every cycle exposed after one
component-union round has no edge-disjoint host cycle on the same
support. The actual endpoint
\href{https://github.com/alejandrozu/openmath-2026-judging/blob/70c5ec3107658182f22c207d4c50f53f9cb3e602/robert-review-2026-10-07/lean/InfiniteAffinePaletteBarrier.lean\#L54}{\texttt{OpenMathReview.InfiniteAffinePaletteBarrier.sidon\_exposed\_cycle\_no\_partner}}
is kernel verified without \texttt{Fintype\ V}. Its proof first derives
the geometric closure from the algebraic Sidon condition, then proves
the actual canonical matching labels, legal swaps, component inducedness
and cycle exclusion. No detector failure or degree bound is a
conclusion-shaped premise. If \(|S|\geq4\), Theorem H simultaneously
supplies an uncoloured pair, emphasizing the limitation of the detector.

\textbf{Theorem R2 (formal two-round sharpness).} For every \(k\geq2\),
the canonical parabola-coloured full host over \(\mathbb F_{3^k}\) fails
the zero/one-round detector, but two legal component-union rounds expose
an actual pair on a common support of exactly 54 vertices. Choose
\(t\notin\{0,1,-1\}\) and the four colours \(0,1,t,1+t\). Their
injective rank-three chart transports the \(p=3\) construction into the
full host. The first round swaps selected components with colours \(0\)
and \(t\); the second swaps the marked component with colours \(1\) and
\(1+t\) after the first move. The proof checks component closure, the
literal successive label swaps, matching identification and the
connected final cycles. The endpoint
\href{https://github.com/alejandrozu/openmath-2026-judging/blob/70c5ec3107658182f22c207d4c50f53f9cb3e602/robert-review-2026-10-07/lean/CycleComposition.lean\#L171}{\texttt{RobertPublishable.TWO.family\_exactly\_two\_component\_union\_rounds}}
combines the lower bound and this full-host witness. Thus two is the
minimum number of component-union rounds for this family. The usual
construction describes three component flips followed by one; the
theorem does not assert minimum individual-flip count.

\textbf{Proof of the two-round upper bound.} In characteristic three
choose \(t\notin\{0,1,-1\}\) and canonical colours \(0,1,t,1+t\). Their
parabola vectors are affinely independent over \(\mathbb F_3\), so the
chart \(f:\mathbb F_3^3\to\mathbb F_q^2\) is injective. Its two shores
have 54 vertices. Identify the selected colours with \(0,A,B,C\).

First swap \(0,B\) on the three components \(a=c\), with \(b\)
arbitrary. These are whole global bichromatic six-cycles: their
successor is translation by the selected vector \(e_b\), of order three.
This produces precisely \(Z,B\) above. Next swap \(A,C\) on the
component with left marks \(m_u=(u,0,-u)\). The identity
\(A(m_u)=C(m_{u+1})\) proves component closure; \(A,C\) were unchanged
by round one. The rank-three splice therefore gives the two Hamilton
cycles on the same 54 vertices. Other colours and every host edge remain
unchanged. The one-round induced-support obstruction rules out zero or
one round, while this witness gives two. The count concerns
component-union rounds, not individual flips \cite{ref11}.

\subsection{Quantitative bipartite
Hamiltonicity}\label{quantitative-bipartite-hamiltonicity}

For a balanced bipartite \(d\)-regular graph, write \(W\) for the
biadjacency matrix. Its nontrivial singular-value condition is
\(s_2(W)\leq(1-\delta)d\); the negative adjacency eigenvalue \(-d\) is
inherent in bipartiteness and must not be inserted into a two-sided
absolute-gap theorem.

\textbf{BM quantitative candidate.} The written dossier proposes an
absolute-constant threshold \(d\geq C_{\rm BM}\delta^{-5}(\log n)^3\)
for Hamiltonicity in this one-sided bipartite setting, with a
conservative power-six route also recorded. The graph has positive
degree and even order \(n\geq4\); the asymptotic proof requires its
stated sufficiently-large-order convention. The complete sampling,
router and rounding argument is linked in the repository. This edition
does not call the main theorem Lean verified without an exact audited
endpoint.

Müyesser \href{https://arxiv.org/html/2609.35766v1}{the cited preprint}
proves a power-six two-sided spectral theorem and explicitly notes a
bipartite one-sided analogue after Corollary 1.4. Bradač and Janzer
\href{https://arxiv.org/html/2605.15043v1}{their preprint} already prove
Hamiltonicity and stronger resilience results for regular bipartite
expanders. Consequently generic bipartite Hamiltonicity is background.
The potentially original delta is the explicit power-five quantitative
ledger and its faithful side/parity implementation.

The written power budget is: the sampled layer degree needs
\(D\gtrsim\delta^{-2}\log n\); the router supports
\(k\asymp\delta^2 qn/(\log n)^2\) terminals; and the perturbation slack
permits reservoir fraction \(q\asymp\delta\). Substituting
\(D\asymp kd/n\) gives the candidate power \(\delta^{-5}\). Each of
those statements has technical premises, including matrix sampling,
bounded endpoint multiplicity, actual disjoint paths and balanced
divisibility. A count of exponents is not a substitute for proving them.

Two dependencies now have complete Lean proofs using standard axioms.
B2.2 assumes a real binary \(N\times N\) biadjacency with row and column
sums \(d>0\), \(N\ge2\), \(s_2(W)\le(1-\delta)d\), \(a>0\),
\(0<\delta<1\), and independent uniform \(m\)-subsets \(U,V\), where
\(1\le m\le N\) and \(p=m/N\):

\[
D=pd\ge32(a+4)\delta^{-2}\log(2N)
\quad\Longrightarrow\quad
\Pr\!\left[s_2(W[U,V])>(1-\delta/2)D\right]\le(2N)^{-a}.
\]

Lemma F assumes a finite nonempty balanced bipartite host \(H\), degrees
in \([(1-\eta)d,d]\), \(0<\gamma\le1\), \(\eta\le\gamma^2/8\),
\(d\ge16/\gamma^2\), and \(|\partial_H Z|\ge\gamma d|Z|\) for every
nonempty \(Z\) with \(2|Z|\le|V(H)|\). It constructs a spanning
\(k\)-regular \(F\subseteq H\) with

\[
k=\left\lfloor\left(1-\eta-\frac{\eta+2/d}{\gamma}\right)d\right\rfloor_{+},
\qquad k\ge(1-3\gamma/8)d-1,
\qquad |\partial_F Z|\ge\frac\gamma2 k|Z|
\quad(2|Z|\le|V(H)|).
\]

The degree interval forces \(\eta\ge0\); the cut range is half the total
vertex set. PartitionTransport preserves the host's vertex labels. These
proofs close sampling and factor dependencies; main BM Hamiltonicity and
Erdős 585 remain unproved. See
\href{https://github.com/alejandrozu/openmath-2026-judging/blob/70c5ec3107658182f22c207d4c50f53f9cb3e602/robert-review-2026-10-07/lean/B22.lean\#L26}{B22
Lean},
\href{https://github.com/alejandrozu/openmath-2026-judging/blob/70c5ec3107658182f22c207d4c50f53f9cb3e602/robert-review-2026-10-07/lean/LemmaF.lean\#L18}{LemmaF
Lean},
\href{https://github.com/alejandrozu/openmath-2026-judging/blob/70c5ec3107658182f22c207d4c50f53f9cb3e602/robert-review-2026-10-07/lean/PartitionTransport.lean\#L54}{PartitionTransport
Lean}.

The new graph-semantic proofs verify Hamilton-cycle remainder transfer
and the balanced-deletion cut loss \(t|S|\), rather than the coarser
\(2t|S|\) bound. They retain the Hamiltonicity/regular-factor premises
explicitly and do not certify the main BM theorem. The isolated written
cherry-packing lemma also needs \(d>0\): with \(d=0\), empty shores and
one requested path satisfy its literal inequality but provide no path.
The new Lean counterexample certifies that defect; positive degree
already holds in the intended application. The factor and
vertex-deletion consequences are attached to the accepted Hamiltonicity
input, not separately counted discoveries. Equal deletions from the two
shores are necessary for a spanning bipartite cycle. In edge-expansion
language, the power-five and power-six routes correspond to
inverse-expansion powers ten and twelve. The avoidance-conditioned
extraction bridge E110 remains open; neither route yields an
unconditional improvement to the general upper bound for \(f(n)\). The
unavailable private exponent-eleven and B2 projects are outside this
manuscript's certified claims.

\textbf{Conditional degree-six consequence.} If every nonempty finite
simple bipartite six-regular graph has a pair, then
\(f(n)=\Theta(n\log\log n)\). This is a consequence of published
regular-subgraph theory, rather than a claim that the hypothesis is
proved. Indeed, a spanning bipartite subgraph retains at least half the
edges of any graph. Once the average degree exceeds a suitable constant
times \(\log\log n\), the regular-subgraph theorem \cite{ref4} supplies a
six-regular subgraph. The assumed cycle-pair statement then gives a pair
in the original graph. The lower order is the known construction
\cite{ref3}, completing the conditional conclusion.

\textbf{Conditional polylogarithmic consequence.} Suppose every
bipartite \(r\)-regular graph on \(N\) vertices with
\(r\ge A(\log(2N))^3\) has a pair. The regular-subgraph extraction
theorem of Chakraborti, Janzer, Methuku and Montgomery {[}12, Theorem
1.5{]}, with average-degree requirement \(Cr\log(n/r)\), then yields
\(f(n)=O(n(\log n)^4)\). Choose \(r=\lceil A(\log(2n))^3\rceil\) for
sufficiently large \(n\) and first take a spanning bipartite subgraph;
the extracted regular graph has order at most \(n\), so the hypothesis
applies. This reduction is conditional. The separate extraction bridge
in the source programme is open, and the completed sampling and factor
lemmas do not supply it.

\section{Proof repository, novelty and
reproducibility}\label{proof-repository-novelty-and-reproducibility}

The
\href{https://github.com/alejandrozu/openmath-2026-judging/tree/70c5ec3107658182f22c207d4c50f53f9cb3e602/robert-review-2026-10-07}{review
repository} provides the frozen author sources, new additive Lean
proofs, complete written arguments, graph certificates,
primary-literature comparisons and a result-to-declaration index. The
manuscript remains short by linking proof infrastructure rather than
reproducing every helper declaration and operational audit.

Table~\ref{tab:proof-index} links the selected statements to their theorem or verification evidence.

{
\begin{longtable}[]{@{}
  >{\raggedright\arraybackslash}p{(\linewidth - 2\tabcolsep) * \real{0.3000}}
  >{\raggedright\arraybackslash}p{(\linewidth - 2\tabcolsep) * \real{0.7000}}@{}}
\caption{Selected theorem and verification evidence index}\label{tab:proof-index}\\
\toprule\noalign{}
\begin{minipage}[b]{\linewidth}\raggedright
Claim
\end{minipage} & \begin{minipage}[b]{\linewidth}\raggedright
Exact existing Lean endpoint or evidence
\end{minipage} \\
\midrule\noalign{}
\endhead
\bottomrule\noalign{}
\endlastfoot
Q &
\href{https://github.com/alejandrozu/openmath-2026-judging/blob/70c5ec3107658182f22c207d4c50f53f9cb3e602/robert-review-2026-10-07/lean/Openmath/Proofs/QB5/Statement.lean\#L42}{\texttt{Erdos585.qb5}} \\
Q1 &
\href{https://github.com/alejandrozu/openmath-2026-judging/blob/70c5ec3107658182f22c207d4c50f53f9cb3e602/robert-review-2026-10-07/lean/Openmath/Proofs/QB4/Statement.lean\#L230}{\texttt{Erdos585.exists\_four\_regular\_avoiding\_of\_bipartite\_six\_regular}} \\
T &
\href{https://github.com/alejandrozu/openmath-2026-judging/blob/70c5ec3107658182f22c207d4c50f53f9cb3e602/robert-review-2026-10-07/lean/TwinDensityExtension.lean\#L132}{\texttt{OpenMathReview.TwinDensityExtension.hasPair\_of\_no\_singleton\_twins\_dense}} \\
Finite lower witnesses &
\href{https://github.com/alejandrozu/openmath-2026-judging/blob/70c5ec3107658182f22c207d4c50f53f9cb3e602/robert-review-2026-10-07/lean/Openmath/Proofs/Final.lean\#L112}{\texttt{Erdos585.thirty\_one\_le\_maxEdges\_eleven}},
\href{https://github.com/alejandrozu/openmath-2026-judging/blob/70c5ec3107658182f22c207d4c50f53f9cb3e602/robert-review-2026-10-07/lean/Openmath/Proofs/Final.lean\#L116}{\texttt{Erdos585.thirty\_six\_le\_maxEdges\_twelve}} \\
H &
\href{https://github.com/alejandrozu/openmath-2026-judging/blob/70c5ec3107658182f22c207d4c50f53f9cb3e602/robert-review-2026-10-07/lean/HaarInfinite.lean\#L110}{\texttt{OpenMathReview.HaarInfinite.hasPair\_iff}} \\
Exact-four-colour decomposition &
\href{https://github.com/alejandrozu/openmath-2026-judging/blob/70c5ec3107658182f22c207d4c50f53f9cb3e602/robert-review-2026-10-07/lean/HaarComponents.lean\#L73}{\texttt{RobertPublishable.HaarComponent.every\_connected\_component}} \\
Composite Haar &
\href{https://github.com/alejandrozu/openmath-2026-judging/blob/70c5ec3107658182f22c207d4c50f53f9cb3e602/robert-review-2026-10-07/lean/Openmath/Proofs/HaarCompositeEven.lean\#L441}{\texttt{Erdos585.HaarCompositeEven.hasTwoEdgeDisjointCyclesSameVertexSet\_all}} \\
R &
\href{https://github.com/alejandrozu/openmath-2026-judging/blob/70c5ec3107658182f22c207d4c50f53f9cb3e602/robert-review-2026-10-07/lean/Openmath/Proofs/ParabolaBarrier104.lean\#L145}{\texttt{Erdos585.ParabolaBarrier104.arbitrary\_polylog\_barrier}} \\
S &
\href{https://github.com/alejandrozu/openmath-2026-judging/blob/70c5ec3107658182f22c207d4c50f53f9cb3e602/robert-review-2026-10-07/lean/R1.lean\#L13}{\texttt{RobertPublishable.SUB.theorem\_R1}} \\
General Sidon barrier &
\href{https://github.com/alejandrozu/openmath-2026-judging/blob/70c5ec3107658182f22c207d4c50f53f9cb3e602/robert-review-2026-10-07/lean/InfiniteAffinePaletteBarrier.lean\#L54}{\texttt{OpenMathReview.InfiniteAffinePaletteBarrier.sidon\_exposed\_cycle\_no\_partner}} \\
S1 &
\href{https://github.com/alejandrozu/openmath-2026-judging/blob/70c5ec3107658182f22c207d4c50f53f9cb3e602/robert-review-2026-10-07/lean/PrescribedVertex.lean\#L89}{\texttt{RobertPublishable.SUB.prescribed\_vertex\_equivalence}} \\
Two-round sharpness &
\href{https://github.com/alejandrozu/openmath-2026-judging/blob/70c5ec3107658182f22c207d4c50f53f9cb3e602/robert-review-2026-10-07/lean/CycleComposition.lean\#L171}{\texttt{RobertPublishable.TWO.family\_exactly\_two\_component\_union\_rounds}} \\
BM & Main Hamiltonicity endpoint remains unverified; exact helper scopes
in the theorem index \\
\end{longtable}
}

The frozen audit used Lean 4.33.1, Mathlib revision
\texttt{0df444a360eaa60ab8c11dca51a86af692955474}, and Formal
Conjectures revision \texttt{137aec5c7abd3aa61f7a73138a97279acfc79e93}.
Selected existing declarations use only \texttt{propext},
\texttt{Classical.choice} and \texttt{Quot.sound}; the finite quartic
case check uses kernel reduction, not \texttt{native\_decide}. New
declarations have their own axiom outputs. A proof replay and its
dependency pin establish what was checked, not who discovered the
mathematics first.

The novelty record distinguishes an exact located predecessor, a result
not implied by that predecessor, and an unresolved priority question.
Our strongest present publication candidates are the precise extraction
threshold, the restricted twin-forcing theorem, the specified
recolouring obstruction and any Haar cases surviving exact decomposition
comparisons. Finite examples and classifications must retain their
evidence and attribution boundaries. Substitution preservation and the
prescribed-vertex equivalence are formally complete; their historical
priority remains unresolved. BM retains its explicit main proof
obligations. Supporting interfaces and classical ingredients do not
become additional discoveries by proximity to compiled code.

\textbf{Acknowledgements and authorship.} Robert Huynh supplied the
original mathematics, constructions and formal development. Alejandro
Zarzuelo Urdiales contributed literature comparison, exposition,
independent verification and the additive formal proofs identified in
the companion. Original source copyrights, Apache-2.0 notices and
library attributions are retained.

\textbf{AI assistance.} The original development used AI agents under
Robert Huynh's direction. Codex assisted the present literature
synthesis, drafting and additive formalization. The formal audit
certifies the selected encoded statements; responsibility for the
mathematical claims, attribution and final text belongs to the human
authors.

\appendix
\section{Additional source reports and the complete
public-result
map}\label{appendix-a.-additional-source-reports-and-the-complete-public-result-map}

\textbf{P2 and near-additive quartic questions.} The author reports an
ordinary proof that every simple graph, not necessarily bipartite, with
\(n\ge2\), maximum degree at most six and at least \(3n-2\) edges has a
nonempty quartic subgraph. Its Tutte-budget/C1 proof has AI reviews but
no selected formal endpoint here; it remains an attributed appendix
claim. The stronger bipartite finite census at \(3n-7\), the \(3n-8\)
witnesses, and the failed general-graph P4 proof attempt are also
distinct. A finite census through a stated order does not prove a
universal additive threshold.

\textbf{The B6 census programme.} The reported W1 counts through
fourteen vertices, bipartite counts through eighteen,
written-plus-computational extension through twenty-three, and resulting
prescribed-vertex B6 claim through twenty-four belong to the
supplementary record. The S6 census at twenty-two is one main method
with a small independent spot-check, not a universal theorem. False
spanning-decomposition and marked-edge routes have explicit
counterexamples; they do not refute B6. Those records can guide research
without being relabelled as fresh full replay or accepted discoveries.

Here W1 is the assertion that a simple graph with maximum degree at most
six and at least \(3n-4\) edges contains a pair. The reported
unrestricted census covers \(2\le n\le14\), and its bipartite
counterpart covers \(2\le n\le18\). The source's additional reductions
and censuses assert the bipartite version through \(n=23\); since a
vertex-deleted six-regular graph has \(3n-3\) edges, this would give the
prescribed-vertex assertion for nonempty bipartite six-regular hosts
through order 24. These are finite source reports, not a proof of the
universal conjecture.

\textbf{The S6 condition.} For a bipartite six-regular graph \(B\),
delete a vertex \(o\) and write \(\Gamma=B-o\). Its six neighbours of
\(o\) are the degree-five ports on the larger shore. Call \(\Gamma\)
sparse when \(g(S)=6|S|-2e_\Gamma(S)\ge10\) for every
\(2\le|S|\le|V(\Gamma)|-1\). S6 asserts that, under this hypothesis,
deleting any port \(y\) from \(\Gamma\) leaves two edge-disjoint
spanning Hamilton cycles. The source relates this sparsity to the E10
cut condition: every cut of \(B\) with at least two vertices on each
side has at least ten edges. Its 22-vertex S6 census and 0.376-percent
spot-check are evidence at that order only; the all-order statement
remains open.

\textbf{Why a spanning-cycle shortcut fails.} Take two copies of
\(K_{5,5}\) with shores \(\{0,\ldots,4\},\{10,\ldots,14\}\) and
\(\{5,\ldots,9\},\{15,\ldots,19\}\). Add
\(0\!-!15,1\!-!16,2\!-!17,3\!-!18,4\!-!19\) and \(10\!-!5\). This is the
source's 20-vertex, 56-edge sparse example. Each spanning Hamilton cycle
must cross between the pieces. The first piece is balanced, so its cycle
uses equally many crossing edges of the two shore orientations; the edge
\(10\!-!5\) is the sole crossing edge of one orientation and must be
used. Two edge-disjoint spanning Hamilton cycles are therefore
impossible, although an internal \(K_{4,4}\) supplies a pair on a
smaller support. The separate marked-edge rule is also false: the source
supplies six 11-vertex, 30-edge certificates where removal of two marked
disjoint edges is pair-free but no pair places those marks on different
cycles. These counterexamples reject the specified shortcuts, not B6.

\subsection{The original 22 headline
rows}\label{the-original-22-headline-rows}

This catalogue preserves source scope rather than awarding one
contribution per row. ``Formal'' refers to the actual matched graph
statement; ``reported'' retains the author's computational/written
coverage and its restrictions. The eight families remain QB, TWIN,
FINITE, REG5, SUB, HAAR, ONE and BM. Most auxiliary census and
obstruction rows lie within those programmes or their open-question
record.

Table~\ref{tab:headline-catalogue} records the original 22 headline claims with their individual evidence boundaries.

{
\begin{longtable}[]{@{}
  >{\raggedleft\arraybackslash}p{(\linewidth - 4\tabcolsep) * \real{0.1300}}
  >{\raggedright\arraybackslash}p{(\linewidth - 4\tabcolsep) * \real{0.4800}}
  >{\raggedright\arraybackslash}p{(\linewidth - 4\tabcolsep) * \real{0.3900}}@{}}
\caption{Catalogue of the original 22 headline claims and evidence boundaries}\label{tab:headline-catalogue}\\
\toprule\noalign{}
\begin{minipage}[b]{\linewidth}\raggedleft
Original row
\end{minipage} & \begin{minipage}[b]{\linewidth}\raggedright
Claim retained
\end{minipage} & \begin{minipage}[b]{\linewidth}\raggedright
Current evidence and boundary
\end{minipage} \\
\midrule\noalign{}
\endhead
\bottomrule\noalign{}
\endlastfoot
1 & QB5: bipartite \(\Delta\le6,\ n\ge3,\ e\ge3n-5\) gives a nonempty
quartic & Formal universal theorem; a quartic is not a pair. \\
2 & QB4 at \(3n-4\), and a quartic in every vertex-deleted bipartite
six-regular host & Formal earlier consequence; no Hamilton splitting
conclusion. \\
3 & P2: general simple \(\Delta\le6,\ n\ge2,\ e\ge3n-2\) & Author's
ordinary reviewed proof; no selected formal P2 endpoint. \\
4 & No-quartic edge census: \(3n-7\) suffices for \(4\le n\le19\);
extrema at 17--19 are 43,46,49 & Source-reported two-method finite
computation; not a universal threshold. \\
5 & Pair-free five-regular graph on 18 vertices & Formal actual witness;
minimum order is a separate census claim. \\
6 & Pair-free five-regular graph on 32 vertices & Formal actual witness,
same REG5 family. \\
7 & Bipartite pair-free five-regular graph on 104 vertices & Formal
first assignment; alternate assignment computational. \\
8 & No pair-free five-regular graph through 16; claimed minimum 18 & One
complete 16-vertex method plus partial second method; qualification
retained. \\
9 & Exact \(f(n)\) through 7 & Formal upper and lower inequalities,
including trivial orders. \\
10 & \(f(8)=19,\ f(9)=23,\ f(10)=27\) & Finite methods independently
replayed in the earlier audit; not universal formulas. \\
11 & \(f(11)=31,\ f(12)=36\) & Lower witnesses formal; exact upper
exclusions remain source-reported. \\
12 & Explicit linear lower bounds and incidence/subdivision families &
Formal constructions with printed parameter restrictions; no asymptotic
improvement. \\
13 & W1 at \(3n-4\) through 14, bipartite through 18 & Reported finite
census with order-dependent replication; no new complete census here. \\
14 & No pair-free six-regular graph through 15 & Reported direct counts
and W1 consequence; finite evidence only. \\
15 & B6 would imply \(f(n)=\Theta(n\log\log n)\) & Conditional
consequence of published regular-subgraph/lower-bound results; B6
remains open. \\
16 & Substitution for \(d\le7\) and prescribed-vertex B6 equivalence &
Full new formal preservation/equivalence; neither proves B6. \\
17 & Bipartite W1 through 23; prescribed-vertex B6 through 24 &
Written-plus-census source claim with shared/partial replication limits;
not a formal universal result. \\
18 & S6 at 22 vertices & Source-reported main census with 0.376\%
independent spot-check; its sparsity hypothesis is essential. \\
19 & Counterexamples to two proposed B6 routes & Particular
spanning/marked-edge rules fail; the target B6 is not refuted. \\
20 & B(polylog), via an open extraction step, would give
\(O(n(\log n)^4)\) & Conditional published-theory reduction; sampling
and Lemma F are now formal, main BM/extraction still incomplete. \\
21 & One-round canonical recolouring barrier & Formal specified
detector; new full-host two-round sharpness does not concern arbitrary
initial colourings or individual-flip minimum. \\
22 & Twin-core forcing at exact \(3v-2\) & Original formal restricted
host; new no-singleton inequality-range extension belongs to the same
family. \\
\end{longtable}
}

\subsection{Supplementary results and supporting
declarations}\label{supplementary-results-and-supporting-declarations}

The complete names and source files are listed in
\href{https://github.com/roberthuynh/erdos-585/blob/1efc324e155ef0b6a7081c5f695c6c825e7debef/findings/supplement/lean/DECLARATIONS.md}{\texttt{DECLARATIONS.md}}.
They are included in the earlier 106-source replay and its
selected-standard-axiom coverage; this is not 56 additional discoveries.
A compact nonoverlapping grouping is:

\begin{itemize}
\tightlist
\item
  \textbf{Sixteen finite/construction declarations} (rows
  1,35--36,39--44,47,49--50,52--55): hub witnesses, Latin incidence,
  matching subdivisions, rim bounds, deletion counts and small-order
  exclusion tools. Counts and lower inequalities remain part of FINITE,
  not new asymptotic claims.
\item
  \textbf{Twenty-two structured-class/cycle declarations} (rows
  7--11,13--14,20--33,51): cubelike witnesses and forest interfaces;
  Haar rank constructions, cycle splices, exact-four threshold,
  composite parity template; and the Wenger no-six-cycle boundary. Known
  subcases retain their prior context. The formal Haar-extraction
  obstruction says that the Wenger host cannot contain an injective
  abelian three-colour Haar subgraph; it does not say the host is
  pair-free.
\item
  \textbf{Ten transfer limitations and split constructions} (rows
  3--6,12,15,17,19,38,48): explicit counterexamples or restrictions on
  Eulerian splitting, cover transfer, state projection and
  degree-raising rules. They invalidate specified transformations, not
  the parent conjecture.
\item
  \textbf{Eight supporting interfaces} (rows 2,16,18,34,37,45--46,56):
  bipartite balance, density core, labelled cuts/incidence lifts,
  matching probes, port completion, witness counts and small-cut
  confinement. Their hypotheses and actual-cycle interpretation remain
  visible; helper count does not determine scientific credit.
\end{itemize}

The Haar composite template and one-/two-round mechanisms are
substantive results already retained in the main text, not lost among
these groups. The linked comprehensive dossier gives the rank-specific
permutations and splice checks, the exact component choices in the two
rounds, full QB5 covering/hub arguments and operational census records.
It remains useful for expert verification; the concise article and this
catalogue do not claim to reproduce every one of those technical proofs.

\begin{thebibliography}{99}
\bibitem{ref1}
N. Alon, S. Friedland and G. Kalai, Regular subgraphs of almost
regular graphs, Journal of Combinatorial Theory, Series B 37 (1984),
79-91.
\href{https://web.math.princeton.edu/~nalon/PDFS/Publications/Regular\%20subgraphs\%20of\%20almost\%20regular\%20graphs.pdf}{Author-hosted
paper}.

\bibitem{ref2}
D. Chakraborti, O. Janzer, A. Methuku and R. Montgomery,
Edge-disjoint cycles with the same vertex set, Advances in Mathematics
469 (2025), 110228. \href{https://arxiv.org/abs/2404.07190}{Primary
preprint}.

\bibitem{ref3}
L. Pyber, V. Rödl and E. Szemerédi, Dense graphs without
3-regular subgraphs, Journal of Combinatorial Theory, Series B 63
(1995), 41-54. The lower-bound application used here is discussed in
\cite{ref2}.

\bibitem{ref4}
O. Janzer and B. Sudakov, Resolution of the Erdős-Sauer problem
on regular subgraphs, Forum of Mathematics, Pi 11 (2023).
\href{https://arxiv.org/abs/2204.12455}{Primary preprint}.

\bibitem{ref5}
G. H. J. Meredith, Regular n-valent n-connected nonHamiltonian
non-n-edge-colorable graphs, Journal of Combinatorial Theory, Series B
14 (1973), 55-60.
\href{https://doi.org/10.1016/S0095-8956(73)80006-1}{Publisher record}.

\bibitem{ref6}
J.-C. Bermond, O. Favaron and M. Mahéo, Hamiltonian
decomposition of Cayley graphs of degree 4, Journal of Combinatorial
Theory, Series B 46 (1989), 142-153.
\href{https://doi.org/10.1016/0095-8956(89)90040-3}{Publisher record}.

\bibitem{ref7}
Y. Huang, M. Tait and R. Won, Sidon sets and 2-caps in
\(\mathbb F_3^n\), Involve 12 (2019), 995-1003.
\href{https://msp.org/involve/2019/12-6/involve-v12-n6-p06-p.pdf}{Published
primary paper}, particularly Theorems 3.2 and 3.4.

\bibitem{ref8}
D. Eppstein, Hamiltonian Cycles in Subdivided Doubles, Ars
Mathematica Contemporanea 26 (4.02) (2026), 1-9; arXiv: 2510.18359v1 (21
October 2025). \href{https://arxiv.org/html/2510.18359v1}{Primary
theorem and proof}.

\bibitem{ref9}
R. C. Read and R. J. Wilson, An Atlas of Graphs, Oxford
University Press (1998), Chapter 5, p.155.
\href{https://academic.oup.com/book/54439}{Publisher},
\href{https://oeis.org/A000088/a000088_14.pdf}{chapter scan}.

\bibitem{ref10}
S. Bonvicini, T. Pisanski and A. Žitnik, All generalized rose
window graphs are hamiltonian, Graphs and Combinatorics 42, article 27
(2026), Proposition 5.1.
\href{https://doi.org/10.1007/s00373-026-03016-w}{Primary publication}.

\bibitem{ref11}
R. Huynh, Erdős 585 public sources and findings, revision
\texttt{1efc324e155ef0b6a7081c5f695c6c825e7debef}.
\href{https://github.com/roberthuynh/erdos-585/blob/1efc324e155ef0b6a7081c5f695c6c825e7debef/findings/FINDINGS.md}{Frozen
findings}. The companion retains the complete source map and its
associated written proof dossiers.

\bibitem{ref12}
D. Chakraborti, O. Janzer, A. Methuku and R. Montgomery,
Regular subgraphs at every density, Transactions of the American
Mathematical Society (2026), accepted/in press.
\href{https://doi.org/10.1090/tran/9694}{Publisher DOI};
\href{https://wrap.warwick.ac.uk/id/eprint/194663/}{institutional
accepted-version record}. The theorem numbering follows the linked
preprint version in the source findings.
\end{thebibliography}

\section*{Mathematical proof supplement}
{[}RH{]} Frozen QB5 Parts I--IV, with their stated corrections, and
\href{https://github.com/roberthuynh/erdos-585/blob/1efc324e155ef0b6a7081c5f695c6c825e7debef/Openmath/Proofs/QB5/Statement.lean\#L42}{Statement};
{[}QB-cover{]}
\href{https://github.com/roberthuynh/erdos-585/blob/1efc324e155ef0b6a7081c5f695c6c825e7debef/Openmath/Proofs/QB5/Cover.lean\#L800}{refined\_cover}
and CoverCheck 1--3; {[}RH{]} HaarRankTwo; {[}RH{]} HaarWitness/Haar;
{[}RH{]} HaarCompositeDefs/HaarCompositeBlue; {[}RH{]}
HaarCompositeEvenDefs/Red/Blue and
\href{https://github.com/roberthuynh/erdos-585/blob/1efc324e155ef0b6a7081c5f695c6c825e7debef/findings/supplement/lean/Openmath/Proofs/HaarCompositeEven.lean\#L441}{all-parameter
endpoint}; {[}RH{]}
\href{https://github.com/alejandrozu/openmath-2026-judging/blob/70c5ec3107658182f22c207d4c50f53f9cb3e602/robert-review-2026-10-07/lean/CycleComposition.lean\#L171}{actual
full-host composition}.

\end{document}
